\documentclass[a4paper,12pt]{article}
\usepackage{amsmath,amsfonts,amssymb} % Math packages
\usepackage{mathtools} % For defining floor and ceiling functions
\usepackage{amsthm}    % Theorems
\usepackage{bm} 	     % Bold math
\usepackage{bbm}	     % More bold math?
\usepackage{array}     % Better tables
\usepackage{physics}   % derivatives and partials
\usepackage{float}     % Better positioning [H]
\usepackage{framed}    % Framed boxes
\usepackage{geometry}  % Required for adjusting page dimensions and margins
\usepackage{graphicx}  % Include images
\usepackage{multirow}  % multicolumn tables
\usepackage{pgfplots}  % Create plots in latex
\usepackage{siunitx}   % SI unit system
\usepackage{listings}  % Code environments
\usepackage[shortlabels]{enumitem} %lets us change the enumeration to (a), (b), ...
\usepackage{booktabs}  % Better looking horizontal rules
\usepackage{makecell}  % Pre-formatting of column heads
\usepackage{hyperref}  % Clickable links
\usepackage{tikz}      % Graphs
\usepackage{cancel}    % Cancel out terms
\usetikzlibrary{shapes,arrows,positioning} % Graph arrows
\usetikzlibrary{automata,positioning,fit,shapes.geometric,backgrounds} % Node grouping
\usetikzlibrary{calc, tikzmark}  % Marking equations
\renewcommand{\arraystretch}{1.6}

\usepackage{mathtools}
\newcommand{\defeq}{\vcentcolon=}
\newcommand{\eqdef}{=\vcentcolon}

\newtheorem*{definition}{Definition}
\newtheorem*{theorem}{Theorem}
\newtheorem*{lemma}{Lemma}
\theoremstyle{definition}
\newtheorem{question}{Question}
\pgfplotsset{width=10cm,compat=1.9}

% stats commands
\newcommand{\E}{\mathbb{E}}
\newcommand{\Var}{\mathbb{V}\mathrm{ar}}
\newcommand{\Cov}{\mathbb{C}\mathrm{ov}}
\newcommand{\Risk}{\mathcal{R}}
\newcommand{\Bias}{\mathrm{Bias}}
\newcommand{\MSE}{\mathrm{MSE}}
\newcommand{\Normal}{\mathcal{N}}
\newcommand{\Exponential}{\mathrm{Exp}}
\newcommand{\Poisson}{\mathrm{Poisson}}
\newcommand{\Bernoulli}{\mathrm{Bernoulli}}
\newcommand{\Binomial}{\mathrm{Binomial}}
\newcommand{\Uniform}{\mathrm{Uniform}}
\newcommand{\Beta}{\mathrm{Beta}}
\newcommand{\GammaDist}{\mathrm{Gamma}}
\DeclareMathOperator*{\argmax}{arg\,max}
\DeclareMathOperator*{\argmin}{arg\,min}

% Common number sets
\newcommand{\N}{\mathbb{N}}
\newcommand{\R}{\mathbb{R}}
\newcommand{\Q}{\mathbb{Q}}
\newcommand{\Z}{\mathbb{Z}}
\newcommand{\C}{\mathbb{C}}

% Cs stuff
\newcommand{\bigO}{\mathcal{O}}

% Indicator function
\newcommand{\Ind}[1]{\mathbbm{1}_{\{#1\}}}

\DeclarePairedDelimiter\ceil{\lceil}{\rceil}
\DeclarePairedDelimiter\floor{\lfloor}{\rfloor}
\DeclarePairedDelimiter\angleb{\langle}{\rangle}

\geometry{
	paper=a4paper, % Paper size, change to letterpaper for US letter size
  margin=1in, % Margin size
	headheight=14pt, % Header height
	footskip=1.5cm, % Space from the bottom margin to the baseline of the footer
	headsep=1.2cm, % Space from the top margin to the baseline of the header
}

\makeatletter
\renewcommand{\maketitle}{
  \begin{center}
    {\large\bfseries \@title \par}
    \vskip 1em
    {\@author}
  \end{center}
}
\makeatother

\let\tss\textsuperscript % superscript macro
\let\oldtextbf\textbf
\renewcommand{\textbf}[1]{\oldtextbf{\boldmath #1}}

% The title of the report
\title{COMP 588 Project Proposal:}
% Write your name in author
\author{Dowoo Kim, Ryan Jackson, David Zhao}
% Course code and date
\date{MATH 324, Winter 2024}

\begin{document}
\maketitle


We preface this proposal by indicating that our work will extend the results of Cai et al. (2025) in their paper ‘On the Value of Cross-Modal Misalignment in Multimodal Representation Learning’. Specifically, we first aim to reproduce their experimental findings, and then stress-test their theoretical framework by exploring the effects of misalignment in a wider variety of settings, including stronger perturbation bias (more on this later).

\subsection*{Background and Motivation}

Multimodal representation learning seeks to learn a shared representation across multiple modalities, such as images and text, so that semantically related samples from different modalities are mapped close together in the shared representation space. This allows for interesting applications such as cross-modal retrieval, out-of-distribution generalization, and zero-shot learning. \\

Currently, one dominant approach to this problem is Multimodal Contrastive Learning (MMCL), exemplified by models such as CLIP. MMCL uses a contrastive loss such as InfoNCE to pull together the representations of paired samples in the shared embedding space (e.g., an image and its corresponding caption), and push apart the representations of unpaired samples. Despite its success, this approach heavily relies on the perfect semantic alignment of the paired data - an image and caption are assumed to be identical views of the same underlying concept. This can be quite problematic as real-world datasets (e.g., LAION-400M) often exhibit significant cross-modal misalignment. For instance, human-annotated captions often omit critical visual details, vary in quality, or even contain irrelevant/misleading information. \\

Historically, researchers have attempted to deal with this misalignment by viewing it in two opposing ways: as a detrimental noise that should be mitigated, or as a beneficial source of diversity that leads to more robust representations. Cai et al. (2025) reconcile these views by formulating the problem using a latent variable model (LVM), which captures the underlying semantic structure of the data while explicitly modeling the misalignment as a latent variable. Their proposed LVM framework

Within their LVM framework, they identify two types of misalignment mechanisms, each differing by their effects on the latent semantic variables. The first, selection bias, arises when the observed text depends only on a non-empty subset of the semantic variables, thereby inducing information loss. For example, a caption might note the color of an object in an image, but omit its texture or its shape. The second, perturbation bias, occurs when the semantic variables influencing the observed text are corrupted or altered, such as misannotating an apple’s color as ‘black’ instead of ‘red’\footnote{One might notice that the misalignment biases we have targeted so far only include text captions. However, these biases are not limited to any particular modality; our focus on textual annotations stems from the observation that human-generated annotations are more susceptible to such biases}. \\.


\end{document} 
